\documentclass[11pt]{article}

\usepackage[margin=1in]{geometry}
\usepackage[T1]{fontenc}
\usepackage[utf8]{inputenc}
\usepackage{lmodern}
\usepackage{amsmath,amssymb}
\usepackage{booktabs}
\usepackage{array}
\usepackage{enumitem}
\usepackage{longtable}
\usepackage{xcolor}
\usepackage{hyperref}
\usepackage{listings}

\hypersetup{
    colorlinks=true,
    linkcolor=blue,
    urlcolor=blue
}

\lstset{
    basicstyle=\ttfamily\small,
    breaklines=true,
    columns=fullflexible,
    keepspaces=true,
    frame=single
}

\title{Dynamic Port Branch}
\author{}
\date{}

\begin{document}
\maketitle

\section{Purpose}

The \texttt{dynamic\_port} branch is the generic image-based dynamic branch without campus-specific semantic colors. It remains the main time-varying benchmark between the clean static branches and the more semantically structured campus branch.

\section{Current Spatial Design}

The branch now samples:
\begin{itemize}[leftmargin=2em]
    \item \textbf{clustered starts},
    \item \textbf{dispersed one-to-one goals}.
\end{itemize}

So the branch no longer acts like a fully scattered dynamic baseline. It now studies what happens when one agent group begins from a local concentration and must spread toward dispersed destinations while the environment changes over time.

\section{Map Preparation}

This branch uses ordinary thresholded image loading rather than campus semantic colors. It also still applies static-density preprocessing before dynamic loop generation. In that sense, it remains the most generic dynamic image branch in the repository.

\section{Why This Branch Still Matters}

The port branch is now useful for isolating the interaction between:
\begin{itemize}[leftmargin=2em]
    \item clustered initial congestion,
    \item dispersed destination pressure,
    \item time-varying obstacle schedules,
    \item classical vs.\ cyclic movement topology.
\end{itemize}

It therefore sits between the synthetic static baseline and the zone-constrained campus experiments.

\end{document}
